% ============================================================================
% EXEMPLO SIMPLES - LAYOUTS NEWTON PAIVA
% Template básico para criar apresentações rapidamente
% ============================================================================
\documentclass[aspectratio=169]{beamer}

% Usa o sistema de layouts Newton Paiva
\usepackage{../styles/newton-paiva-layouts}

% Pacotes úteis
\usepackage[utf8]{inputenc}
\usepackage[portuguese]{babel}

% ============================================================================
% CONFIGURAÇÕES DA SUA APRESENTAÇÃO
% ============================================================================

% Título e subtítulo
\title{Sua Apresentação Aqui}
\subtitle{Subtítulo Opcional}

% Informações da Newton Paiva
\disciplina{Sua Disciplina}
\instituicao{Centro Universitário Newton Paiva}
\professor{Seu Nome}

% Data
\date{\today}

% Seção atual (para layouts completos)
\secaoatual{Introdução}

% Imagem da capa (opcional - remova se não usar)
% \imagemcapa{assets/images/sua-imagem.png}

% ============================================================================
% INÍCIO DO DOCUMENTO
% ============================================================================
\begin{document}

% ============================================================================
% PÁGINA DE TÍTULO
% ============================================================================

% Escolha um dos três estilos de title page:
% \titlepagepadrao      - Layout padrão (2 colunas com imagem)
% \titlepagemoderno     - Layout moderno (horizontal)
% \titlepageclassico    - Layout clássico (centralizado)

\titlepagepadrao
\begin{frame}[plain]
    \maketitle
\end{frame}

% ============================================================================
% SLIDES DE CONTEÚDO
% ============================================================================

% Escolha um layout para seus slides:
% \layoutpadrao         - Para slides normais
% \layoutapresentacao   - Para apresentações formais
% \layoutaula          - Para aulas
% \layoutexercicio     - Para exercícios
% \layoutsecao         - Para slides de seção
% \layoutlimpo         - Para slides especiais

\layoutpadrao

\begin{frame}{Agenda}
    \begin{block}{Tópicos da Apresentação}
        \begin{itemize}
            \item Primeiro tópico
            \item Segundo tópico
            \item Terceiro tópico
            \item Conclusões
        \end{itemize}
    \end{block}
\end{frame}

\begin{frame}{Primeiro Tópico}
    \begin{columns}
        \begin{column}{0.5\textwidth}
            \begin{block}{Conceitos}
                \begin{itemize}
                    \item Conceito A
                    \item Conceito B
                    \item Conceito C
                \end{itemize}
            \end{block}
        \end{column}
        \begin{column}{0.5\textwidth}
            \begin{exampleblock}{Exemplos}
                \begin{itemize}
                    \item Exemplo 1
                    \item Exemplo 2
                    \item Exemplo 3
                \end{itemize}
            \end{exampleblock}
        \end{column}
    \end{columns}
\end{frame}

% Mudando para layout de seção
\layoutsecao

\begin{frame}
    \begin{center}
        \Huge{\textcolor{newtanorange}{\textbf{Nova Seção}}}
        \vspace{1cm}
        
        \Large{\textcolor{white}{Segunda Parte}}
        
        \vspace{2cm}
        \textcolor{white}{\rule{10cm}{2pt}}
    \end{center}
\end{frame}

% Voltando para layout normal
\layoutpadrao
\secaoatual{Segunda Parte}

\begin{frame}{Segundo Tópico}
    \begin{alertblock}{Importante}
        Este é um ponto importante que merece destaque.
    \end{alertblock}
    
    \begin{block}{Detalhes}
        Aqui você pode colocar os detalhes do seu segundo tópico.
        \begin{enumerate}
            \item Primeiro item
            \item Segundo item
            \item Terceiro item
        \end{enumerate}
    \end{block}
\end{frame}

% Layout para exercícios
\layoutexercicio

\begin{frame}{Exercício Prático}
    \begin{exampleblock}{Questão}
        \textbf{Pergunta:} Como você aplicaria este conceito na prática?
    \end{exampleblock}
    
    \begin{block}{Opções}
        \textbf{a)} Primeira opção\\
        \textbf{b)} Segunda opção\\
        \textbf{c)} Terceira opção\\
        \textbf{d)} Quarta opção
    \end{block}
\end{frame}

% Layout limpo para conclusão
\layoutlimpo

\begin{frame}
    \begin{center}
        \vfill
        \Huge{\textcolor{newtanorange}{\textbf{Conclusões}}}
        \vspace{1cm}
        
        \Large{\textcolor{white}{Principais pontos abordados}}
        \vspace{1cm}
        
        \textcolor{white}{\rule{12cm}{2pt}}
        \vspace{1cm}
        
        \begin{itemize}
            \item Ponto principal 1
            \item Ponto principal 2
            \item Ponto principal 3
        \end{itemize}
        \vfill
    \end{center}
\end{frame}

\begin{frame}
    \begin{center}
        \vfill
        \Huge{\textcolor{newtanorange}{\textbf{Obrigado!}}}
        \vspace{1cm}
        
        \newtonlogo[3cm]
        
        \vspace{1cm}
        \Large{\textcolor{white}{Perguntas?}}
        
        \vspace{1.5cm}
        \textcolor{newtanorange}{\textbf{Seu Nome}}
        \\[0.3cm]
        \textcolor{white}{Centro Universitário Newton Paiva}
        \\[0.3cm]
        \textcolor{white}{\small seu.email@newtonpaiva.br}
        \vfill
    \end{center}
\end{frame}

\end{document}
